\subsection{赋值环的例子}

\begin{enumerate}
\def\labelenumi{(\arabic{enumi})}
\item
  每个域都是赋值环。
\item
  若 \(\mathbf{V}'\) 是域 \(\mathbf{K}'\) 的赋值环，\(\mathbf{K}\) 是 \(\mathbf{K}', \mathbf{V}' \cap \mathbf{K}\) 的子域，则由第 2 节定理 1 (c)，后者是 \(\mathbf{K}\) 的赋值环。
\item
  下面的命题给出赋值环的大量例子：
\end{enumerate}

命题 2. 设 A 是局部环，其极大理想是主理想 Ap。若 \(\bigcap_{n=1}^{\infty} A p^n = (0)\)（例如当 A 是 Noether 时，参见第三章 §3 第 2 节

命题 5 的推论），则 \(\mathbf{A}\) 的理想只有 (0) 与诸 \(\mathbf{A}p^n\)；于是或者 \(p\) 是幂零的，或者 \(\mathbf{A}\) 是赋值环。

设 A 被诸 \(\mathbf{A}p^n\) 滤过，用 \(v\) 表示相应的阶函数（第三章 § 2 第 2 节）。由于

\[
\bigcap_ {n = 1} ^ {\infty} A p ^ {n} = (0),
\]

关系 \(v(x) = +\infty\) 蕴涵 \(x = 0\)。设 \(\mathfrak{a}\) 是一个理想 \(\neq (0)\)，属于 \(\mathbf{A}\)，\(a\) 是 \(\mathfrak{a}\) 中使 \(v\) 取到最小值的一个元素；记 \(v(a) = s\)（\(s \neq +\infty\)）。则 \(\mathfrak{a} \subset \mathrm{Ap}^s\)。特别地，存在 \(u \in \mathbf{A}\) 使 \(a = up^s\)；由于 \(a \notin \mathrm{Ap}^{s+1}\)，有 \(u \notin \mathrm{Ap}\)；于是 \(u\) 可逆且 \(p^s \in \mathbf{A}a \subset \mathfrak{a}\)。由此得出 \(\mathfrak{a} = \mathrm{Ap}^s\)，即得我们的第一个断言。同样可见，每个元素 \(a \neq 0\)（\(\mathbf{A}\) 的）都可写成 \(a = up^{v(a)}\) 的形式，其中 \(u\) 可逆。若 \(a' = u'p^{v(a')}\)（\(u'\) 可逆）是 \(\mathbf{A}\) 的另一个非零元素，则 \(aa' = uu'p^{v(a)+v(a')}\)；于是，若 \(p\) 不是幂零的，则 \(aa' \neq 0\)，从而 \(\mathbf{A}\) 是整环。这时，由于 \(\mathbf{A}\) 的理想所成之集按包含关系是全序的，我们断定 \(\mathbf{A}\) 是赋值环（定理 1 (e)）。

例如，若 \(p\) 是素数，则局部环 \(\mathbf{Z}_{(p)}\) 是赋值环。设 \(\mathbf{B} = \mathbf{K}[\mathbf{X}_1, \ldots, \mathbf{X}_n]\) 是 \(n\) 个未定元在域 \(\mathbf{K}\) 上的多项式环；理想 \(\mathbf{BX}_1\) 是素的，因为 \(\mathbf{B} / \mathbf{BX}_1\) 同构于 \(\mathbf{K}[\mathbf{X}_2, \ldots, \mathbf{X}_n]\)；于是 \(\mathbf{B}_{\mathbf{BX}_1}\) 是赋值环；它由有理函数 \(\mathbf{PQ}^{-1}\) 组成，其中 \(\mathbf{P}\)、\(\mathbf{Q}\) 是多项式，且 \(\mathbf{Q}(0, \mathbf{X}_2, \ldots, \mathbf{X}_n) \neq 0\)。

* 更一般地，我们将看到，若 F 是

\[
\mathrm{B} = \mathrm{K} [ \mathrm{X} _ {1}, \dots , \mathrm{X} _ {n} ],
\]

的一个极端元，则 \(B_{BF}\) 是赋值环（参见第七章 § 3 第 5 节）。*

域 K 上一个未定元的形式幂级数环 K{[}{[}T{]}{]} 是 Noether 局部整环，其极大理想是主的；因此它是赋值环。另一方面，两个未定元的形式幂级数环 K{[}{[}T₁, T₂{]}{]} 虽是 Noether 局部整环，却不是赋值环，因为元素 T₁、T₂ 中没有一个是另一个的倍数。

命题 3. 设 A 是主理想整环，K 是它的分式域。K 的包含 A 且异于 K 的赋值环恰是形如 \(A_{Ap}\) 的环，其中 p 是 A 的极端元。

显然 \(A_{Ap}\)（p 为极端元）是包含 A 且异于 K 的赋值环（命题 2）。反之，设 V 是异于 K 且包含 A 的赋值环。由于 \(V \neq K\)，\(m(V)\) 包含一个元素 \(x \neq 0\)；记 x = a/b，其中 \(a \in A\)、\(b \in A\)，可见 \(A \cap m(V)\) 包含非零元素 a。由于 \(A \cap m(V)\) 是素的，它形如 Ap，其中 p 在 A 中是极端的。于是 \(A_{Ap} \subset V\)，\(pA_{Ap} \subset m(V)\)，从而 V 支配 \(A_{Ap}\)；由于 \(A_{Ap}\) 是 K 的赋值环（命题 2），有 \(V = A_{Ap}\)。

推论 1. 域 \(\mathbf{Q}\) 的每个异于 \(\mathbf{Q}\) 的赋值环都形如 \(\mathbf{Z}_{(p)}\)，其中 \(p\) 是素数。

Q 的每个子环都包含 Z。

推论 2. 设 K 是域，K(X) 是 K 上一个未定元的有理函数域，V 是 K(X) 的包含 K 且异于 K(X) 的赋值环。若 X ∈ V，则存在不可约多项式 P ∈ K{[}X{]}，使 \(V = (K{[}X{]})_{(P)}\)；否则 V 是 K{[}X⁻¹{]} 在素理想 X⁻¹K{[}X⁻¹{]} 处的局部环（换言之，V 由分式 A/B 组成，其中 A ∈ K{[}X{]}、B ∈ K{[}X{]}，且 deg(A) ≤ deg(B)）。

若 \(\mathbf{X} \in \mathbf{V}\)，则 \(\mathbf{K}[\mathbf{X}] \subset \mathbf{V}\)，所述断言由命题 3 得出。若 \(\mathbf{X} \notin \mathbf{V}\)，则 \(\mathbf{X}^{-1} \in \mathbf{V}\)，从而 \(\mathbf{K}[\mathbf{X}^{-1}] \subset \mathbf{V}\)；于是 \(\mathbf{V}\) 是 \(\mathbf{K}[\mathbf{X}^{-1}]\) 在某个素理想 \(p\) 处的局部环（命题 3），且这个理想包含 \(\mathbf{X}^{-1}\)，因为 \(\mathbf{X}^{-1}\) 在 V 中不可逆；于是 \(p = X^{-1}K[X^{-1}]\)，因为后一理想是极大的。最后考虑有理函数 \(A(X)/B(X)\)，其中 A、B 是次数分别为 a 与 b 的多项式；则 \(A(X) = X^{a}A'(X^{-1})\)，\(B(X) = X^{b}B'(X^{-1})\)，其中 \(A'\)、\(B'\) 是满足 \(A'(0) \neq 0\) 与 \(B'(0) \neq 0\) 的多项式；于是，为使 \(A(X)/B(X)\) 属于 \(K[X^{-1}]\) 在 \(X^{-1}K[X^{-1}]\) 处的局部环，必须且只需 \(a \leqslant b\)。

